\documentclass[10pt,a4paper]{amsart}
\usepackage[margin=19mm]{geometry}
\usepackage{amsmath,amssymb,mathtools}
\usepackage[scr=rsfs,cal=euler]{mathalfa}
\usepackage{xfrac}
\usepackage[unicode,hidelinks,bookmarks=false]{hyperref}
\newcommand{\ie}{i.e.\ }
\newcommand{\cf}{cf.\ }
\newcommand{\resp}{respectively\ }
\newcommand{\Subsection}{\subsection}
\providecommand{\nobreakdash}{\nobreak\hbox{-}}
\setlength{\emergencystretch}{2em}
\multlinegap0pt
\usepackage{amssymb}
\usepackage{xcolor}
\newtheorem{lemma}{Lemma}
\DeclareMathOperator{\Ber}{Ber}
\newcommand{\E}{\mathbb{E}}
\renewcommand{\P}{\mathbb{P}}
\DeclareMathOperator{\Var}{Var}
\newcommand{\beq}{\begin{eqnarray*}}
\newcommand{\beqn}{\begin{eqnarray}}
\newcommand{\eeq}{\end{eqnarray*}}
\newcommand{\eeqn}{\end{eqnarray}}
\renewcommand{\vec}[1]{\textup{{\bf #1}}}
\title{TV homogenization inequalities}
\author{Aryeh Kontorovich}
\date{Source: arXiv:2601.04079v3 (CC BY 4.0)}
\begin{document}
\maketitle

\noindent\emph{Source and licence.} TV homogenization inequalities. Version arXiv:2601.04079v3 (CC BY 4.0), distributed by arXiv under the Creative Commons Attribution 4.0 International licence, http://creativecommons.org/licenses/by/4.0/, as recorded in the arXiv licence metadata for this exact version and shown on the abstract page https://arxiv.org/abs/2601.04079v3. Full licence notice: \texttt{licenses/arxiv-2601.04079-CC-BY-4.0.txt}.\par\smallskip

\noindent\textit{Editorial scope.} Counted unit: the article's lemma lem:J-global (source0.tex lines 1040--1048). The article's complete original proof of it is delivered with it (source0.tex lines 1050--1172, one \texttt{proof} environment of 123 lines). No mathematics is written, repaired or re-derived: the statement and the proof are the article's own lines, in the article's own notation, and no retained line is substituted or re-flowed.  No further source-local context is needed: every cross-reference of every retained line resolves inside the two delivered spans.  Nothing was omitted from any retained span, so the package carries no omission record.  No retained line contains a citation, so the wrapper carries no bibliography and no bibliography entry.  No retained line contains a figure environment, an image-inclusion command or drawing code, so no image is delivered and the package declares no asset dependency.  Editorial formatting only, in addition: an AMS \texttt{amsart} preamble that restates the article's own declarations and macros used by the retained lines, namely the environments \texttt{lemma} on separate counters, together with the standard packages those lines need.  The retained environments print in this document as Lemma 1 in order of appearance, whereas the article prints the same items with the numbering of its own full document.  The complete native source of the article as distributed in the arXiv e-print for this exact version is bundled unchanged as \texttt{sources/192266/source0.tex}.
\begin{lemma}
\label{lem:J-global}
Let $X,Y,\sigma_{\vec{p}}^2=\Var(X),g,\mu_{\vec{p}},\Delta,J$ be defined as above.
Then
\[
  J \;\le\;
2\,\Delta\,\bigl(\sigma_{\vec{p}}^2+1+\Delta^2\bigr).
\]
\end{lemma}

\begin{proof}
For $i\in[n]$,
write $\Delta_i:=p_i-q_i\ge0$ with $\sum_i\Delta_i=\Delta$. For
$t\in[0,1]$ set
\[
  r_i(t) := q_i + t\Delta_i,\quad
  X_i(t)\sim\Ber(r_i(t))\ \text{independent},\quad
  S(t):=\sum_{i=1}^n X_i(t),\quad
  T_i(t):=S(t)-X_i(t).
\]
Then $S(0)=Y$ and $S(1)=X$ in distribution. For fixed $n$, each
$S(t)$ has a distribution that is a finite polynomial in the parameters
$\{r_i(t)\}$, so for each $k$ and $i$ the function
\beq
F_k(t) &:=&
\P(S(t)\ge k) \\
&=&
r_i(t)\,\P(T_i(t)\ge k-1)\;+\;(1-r_i(t))\,\P(T_i(t)\ge k)
.
\eeq
is a polynomial in $t$ and hence differentiable on $[0,1]$.
Differentiating,
\[
  \frac{\partial}{\partial r_i}F_k
  = \P(T_i(t)\ge k-1)-\P(T_i(t)\ge k)
  = \P(T_i(t)=k-1).
\]
By the chain rule and the fact that $r_i'(t)=\Delta_i$,
\[
  \frac{\mathrm{d}}{\mathrm{d}t}\,\P(S(t)\ge k)
  = \sum_{i=1}^n \Delta_i\,\P\bigl(T_i(t)=k-1\bigr),
  \qquad k\in\mathbb Z.
\]
Integrating from $0$ to $1$ and using $S(0)=Y$,
$S(1)=X$ in distribution,
\[
  g(k)
  = \P(X\ge k)-\P(Y\ge k)
  = \int_0^1 \sum_{i=1}^n \Delta_i\,\P\bigl(T_i(t)=k-1\bigr)\,\mathrm{d}t,
  \qquad k\in\mathbb Z.
\]
In this formulation, the functional $J=\sum_k (k-\mu_{\vec{p}})^2 g(k)$
becomes
\begin{align*}
  J
  &= \int_0^1 \sum_{i=1}^n \Delta_i
       \sum_{k\in\mathbb Z} (k-\mu_{\vec{p}})^2 \P\bigl(T_i(t)=k-1\bigr)\,\mathrm{d}t.
\end{align*}
Changing variable $j=k-1$,
\[
  \sum_{k\in\mathbb Z} (k-\mu_{\vec{p}})^2 \P\bigl(T_i(t)=k-1\bigr)
  = \E_t\bigl[(T_i(t)+1-\mu_{\vec{p}})^2\bigr],
\]
where $\E_t$ denotes expectation under the product law with
parameters $\{r_j(t)\}_j$.
Hence
\begin{equation}\label{eq:J-Et}
  J
  = \int_0^1 \sum_{i=1}^n \Delta_i\,
       \E_t\bigl[(T_i(t)+1-\mu_{\vec{p}})^2\bigr]\,\mathrm{d}t.
\end{equation}
We bound $\E_t[(T_i(t)+1-\mu_{\vec{p}})^2]$ uniformly in $t,i$
via the decomposition
$\E[(Z-\alpha)]^2=\Var(Z)+(\E Z-\alpha)^2$,
where $Z=T_i(t)+1$ and $\alpha=\mu_{\vec{p}}$.
To bound the variance term, note that
$f(x)=x(1-x)$
is $1$-Lipschitz on $[0,1]$, and so
\[
  |f(r_j(t))-f(p_j)| \le |r_j(t)-p_j| = (1-t)\Delta_j.
\]
Hence
\[
  \Var_t(S(t))
  = \sum_j f(r_j(t))
  \le \sum_j f(p_j) + (1-t)\sum_j\Delta_j
  = \sigma_{\vec{p}}^2 + (1-t)\Delta
  \le \sigma_{\vec{p}}^2 + \Delta
  \le\sigma_{\vec{p}}^2 + 1 + \Delta^2.
\]
Since $T_i(t)=S(t)-X_i(t)$ with $X_i(t)$ Bernoulli,
\beqn
\label{eq:VarTi-bound}
  \Var_t(T_i(t))
  \le \Var_t(S(t))
  \le \sigma_{\vec{p}}^2 + 1 + \Delta^2.
\eeqn

To bound the bias term, we have
\[
  \E_t[S(t)] = \mu_{\vec{q}} + t\Delta,
\]
and $\mu_{\vec{p}}=\mu_{\vec{q}}+\Delta$, so $\E_t[S(t)]-\mu_{\vec{p}}=-(1-t)\Delta$.
Also $r_i(t)=q_i+t\Delta_i$, so
\[
  \E_t[T_i(t)]
  = \mu_{\vec{q}} + t\Delta - (q_i + t\Delta_i),
\]
and hence
\[
  M_i(t) := \E_t[T_i(t)+1-\mu_{\vec{p}}]
  = 1 - q_i - \Delta + t(\Delta-\Delta_i).
\]
It is straightforward to see that
$|M_i(t)| \le \max(1,\Delta)$
and thus,
\beq
  \E_t\bigl[(T_i(t)+1-\mu_{\vec{p}})^2\bigr]
  &=& \Var_t(T_i(t)) + M_i(t)^2
  \\
  &\le&
 (\sigma_{\vec{p}}^2+1+\Delta^2) + (\sigma_{\vec{p}}^2+1+\Delta^2)
  = 2(\sigma_{\vec{p}}^2+1+\Delta^2).
\eeq
Substituting into \eqref{eq:J-Et} gives
\[
  J
  \le \int_0^1 \sum_i \Delta_i\,
         2(\sigma_{\vec{p}}^2+1+\Delta^2)\,\mathrm{d}t
  = 2\,\Delta\,(\sigma_{\vec{p}}^2+1+\Delta^2),
\]
as claimed.
\end{proof}

\end{document}
